% Standalone research appendix fragment. Not included in paper.tex here.
\section{An obstruction to the aligned doubled-Hadamard involution model}
\label{app:c2-norm-obstruction}

The statement concerns a specified geometric involution, not an arbitrary
combinatorial automorphism of an incidence design. In particular it does not
exclude the disjoint union of two differently labelled Hadamard designs.

\begin{theorem}\label{thm:c2-norm-obstruction}
Over a field of characteristic different from two, there is no realization
of the following sign-twisted aligned doubled-Hadamard model. The eight
polynomials are
\[
 F_i^\pm(T)=\pm A_i(T^2)+T B_i(T^2),\qquad
 \deg A_i=3,\quad\deg B_i\le2,\qquad 0\le i\le3.
\]
Their eight signed leading coefficients are distinct. Six distinct nonzero
values $e_{ij}$ give the old nodes $\pm\sqrt{e_{ij}}$, at which all four
polynomials with labels $i,j$ agree. The other sixteen nodes are the roots
of four monic quadratics and their reversals, indexed by the even transversal
sign patterns modulo simultaneous sign reversal; the four selected
polynomials agree at each corresponding node. All twenty-eight nodes are
distinct.
\end{theorem}

\begin{proof}
Subtract a common odd polynomial and rescale to arrange $B_0=0$ and $a_0=1$,
where $a_i$ is the leading coefficient of $A_i$. Put
\[
 C_i(Y)=\prod_{j\ne i}(Y-e_{ij}),\qquad A_i=a_iC_i.
\]
The all-plus quadratic is $L_0=T^2-\sigma T+p$. Its roots and their negatives
are distinct and nonzero, so $\sigma p\ne0$. Write
$R_0(Y)=(Y+p)^2-\sigma^2Y$. All $R_0(e_{ij})$ are nonzero.

The difference $F_i^+-F_0^+$ is divisible by $(T^2-e_{0i})L_0$.
Write its quadratic quotient as
$(a_i-1)T^2+\beta_iT+z_i$. Comparing parity gives, without any division by
$p+e_{ij}$ or $C_i(-p)$,
\begin{equation}\label{eq:c2-beta-identity}
 (Y+p)B_i=\sigma(C_0-a_iC_i)+R_0\beta_i(Y-e_{0i}).
\end{equation}
At $e_{ij}$, subtracting these identities yields
\begin{equation}\label{eq:c2-beta-edge}
 \beta_i(e_{ij}-e_{0i})=\beta_j(e_{ij}-e_{0j}).
\end{equation}
Thus the three $\beta_i$ either all vanish or are all nonzero and pairwise
distinct. The all-zero case makes the residual quadratic of $F_i^+-F_0^+$
even, identifying its prescribed transversal with its reversal. This
contradicts distinctness of the domain. Hence every $\beta_i$ is nonzero
and the three are distinct.

Let $u_i=[Y^2]B_i$. Equation~\eqref{eq:c2-beta-identity} gives
$u_i=\beta_i-\sigma(a_i-1)$. For a general pair the monic quadratic
remaining after removal of its old edge and $L_0$ has linear coefficient
$(\beta_i-\beta_j)/(a_i-a_j)$, where $\beta_0=0$. Complementary pairs have
reversed residual quadratics. Setting $d_i=a_i-1$, their three slope
conditions are
\[
 \begin{pmatrix}
 \beta_2-\beta_3&\beta_1&-\beta_1\\
 \beta_2&\beta_1-\beta_3&-\beta_2\\
 \beta_3&-\beta_3&\beta_1-\beta_2
 \end{pmatrix}
 \begin{pmatrix}d_1\\d_2\\d_3\end{pmatrix}=0.
\]
Put $S=\beta_1+\beta_2+\beta_3$. This matrix kills the vector
$v_i=\beta_i(S-2\beta_i)$, and its two by two minors include, up to sign,
these three entries. They cannot all vanish: otherwise $S=2\beta_i$ for
all $i$, which on summing forces $S=0$ and then $\beta_i=0$.
The matrix has rank two and therefore
\begin{equation}\label{eq:c2-amplitudes}
 a_i=1+t\beta_i(S-2\beta_i),\qquad
 \frac{u_i}{\beta_i}=1-\sigma t(S-2\beta_i),\qquad t\ne0.
\end{equation}
In particular the last expression is a nonconstant affine function of
three distinct arguments.

At a new node with received value $V$, squaring the selected equalities
gives $(V-TB_i)^2=A_i^2$ and $V^2=A_0^2$. Consequently
\[
 H_{ij}=YB_iB_j(B_i-B_j)-B_j(A_i^2-A_0^2)+B_i(A_j^2-A_0^2)
\]
vanishes at all eight new square values. It also vanishes at the three
old values $e_{0i},e_{0j},e_{ij}$. Since $\deg H_{ij}\le8$, it vanishes
identically. Its leading coefficient implies
$u_i(a_j^2-1)=u_j(a_i^2-1)$. All $a_i^2-1$ are nonzero; hence either all
$u_i$ vanish or none do. The former is incompatible with
\eqref{eq:c2-amplitudes}. Every $B_i$ therefore has degree two.

We claim $\gcd(B_1,B_2,B_3)=1$. Each $B_i$ vanishes at the distinct point
$e_{0i}$, so a common divisor has degree at most one. A common root $r$
would give $B_i=u_i(Y-r)(Y-e_{0i})$. At every $e_{ij}\ne r$, equality of
the $B_i$ and \eqref{eq:c2-beta-edge} then give
$u_i/\beta_i=u_j/\beta_j$. At least two of the edges $12,13,23$ differ
from $r$, and they connect all three labels. This contradicts
\eqref{eq:c2-amplitudes}.

The identities $H_{ij}=0$ imply that
\[
 C=\frac{A_i^2-A_0^2-YB_i^2}{B_i}
\]
is independent of $i$. Its reduced denominator divides every $B_i$, so
$C$ is polynomial, of degree exactly four. Thus
\[
 R=C^2-4YA_0^2,\qquad
 Q_i(T)=C(T^2)+2T^2B_i(T^2)+2TA_i(T^2)
\]
satisfy $\deg R=8$ and $Q_i(T)Q_i(-T)=R(T^2)$.
At every new square value some $B_i$ is nonzero, and the received-value
equation gives $C=-2TV$. Hence $R$ has precisely the eight distinct
nonzero new square values as its roots. They are disjoint from the old
edge values.

Divide all $Q_i$ by their common leading coefficient. Each now chooses
one root from each of the eight opposite root pairs of $R(T^2)$.
Their constants agree, so pairwise selector distances are even.
Distances zero and eight violate the signed leading-coefficient guards.
Distances two and six are also impossible. Indeed write
$Q_i=ED$, $Q_j=ED(-T)$, with one of $E,D$ quadratic. At an old edge write
$E=e(Y)+To(Y)$ and $D=d(Y)+Tp(Y)$. The equations $A_i=A_j=0$ imply
$ep=od=0$. Since this edge is nonzero and off the norm roots, either
$o=p=0$ or $e=d=0$. The second alternative gives opposite nonzero even
parts $Yop,-Yop$, contrary to $B_i=B_j$. The first would make the odd
part of a quadratic vanish; its two roots would then be opposite,
contrary to the distinct root-square guard. All pairwise distances are
therefore four.

It remains to prove that three such selectors force an edge collision.
Normalize the first selector to all plus. The next two partition the
columns into four cells of size two, with group signs
\[
 (+,+,+,+),\quad (+,+,-,-),\quad (+,-,+,-).
\]
Write the roots in cell $g$ as $m_g\pm n_g$ and set $h_g=m_g^2-n_g^2$.
For a pair of cells $u,v$ with signs $\epsilon_u,\epsilon_v$, the first
and third elementary root sums of their quartic are
\[
 e_1=2(\epsilon_um_u+\epsilon_vm_v),\qquad
 e_3=2(\epsilon_um_uh_v+\epsilon_vm_vh_u).
\]
The preceding odd-part argument gives the three linear conditions
$e_1d_3-e_3d_1=0$ on $h$. Their matrix kills both
$\mathbf1$ and $(m_0,-m_1,-m_2,m_3)$.

Here is an explicit rank certificate, avoiding any generic-rank
assumption. Put
\[
 A=m_0m_1m_2,\ B=m_0m_1m_3,\ C'=m_0m_2m_3,\ D=m_1m_2m_3,
\]
and $S_0=A+B+C'+D$, $S_1=A+B-C'-D$, $S_2=A-B+C'-D$.
For the three pairs of matrix rows, their two by two minors are, up to
sign, respectively $S_0,S_1,S_2$ times
\[
 m_2+m_3,\ m_1+m_3,\ m_1-m_2,\ m_0-m_3,\ m_0+m_2,\ m_0+m_1.
\]
These six factors are nonzero by the three signed leading guards.
Rank at most one would force $(A,B,C',D)=(A,-A,-A,A)$.
If $A\ne0$, this gives $m=(z,-z,-z,z)$ and all three leading sums
vanish. If $A=0$, at most two $m_g$ are nonzero, and two of the three
leading sums coincide up to sign. Both alternatives are forbidden.
Thus the matrix has rank two and
\[
 h=H\mathbf1+K(m_0,-m_1,-m_2,m_3).
\]
For each of the three pairs, substitution into the displayed elementary
sums gives $e_3=He_1$. Since $e_1\ne0$, all three old edge squares are
$-e_3/e_1=-H$, contradicting their distinctness.
\end{proof}

\paragraph{Certificate and scope.}
The determinant identities are integer polynomial identities, reproduced
by \texttt{universal\_beta\_gate.py} and
\texttt{three\_selector\_gate.py} in the research directory
\texttt{hadamard\_locator\_family}. They require no parameter sampling or
Gr\"obner-basis assumption. The proof includes affine opposite-edge
involutions, finite-pole charts, $p+e_{ij}=0$, and coefficient-rank
exceptions, because no involution parametrization is used. It does not
show that a combinatorial design automorphism must lift to a projective
transformation, and makes no assertion about unsymmetrized eight-sextic
banks or disjoint Hadamard designs.
